\chapter{Dodecaphase--Witt intertwining and the exact reading of \texorpdfstring{\(3/11\)}{3/11}}
\label{chap:entrelazador-accion}
\index{intertwiner!dodecaphase--Witt}
\index{Witt!incidence module}
\index{survival!rank-three projector}

The repeated appearance of quotient \(3/11\) on the dodecaphasic screen,
in the polarity of a Witt star, and in nonadic filtration raises
a question stronger than the arithmetic equality of three fractions. The
correct question is operatorial: to determine whether those three appearances are
traces of the same positive sector viewed in different realizations and, if they
are, to construct the maps intertwining them without introducing a subsequent
metrological datum.

The answer has four levels. First, point--hexad incidence
constructs a canonical isometry between the eleven-dimensional dodecaphase
module and an eleven-dimensional eigenspace of the Witt design. Second, the
Witt star acts on that same module, but does so as a signed
involution; its \(3/11\) is an index, not the trace of a positive projector.
Third, the oriented survival sector appears in \(t=7\), where a
unique space of eleven branches contains a rank-three sector. Fourth, the
orientation line of that triplet, together with the internal witnesses
\(B_K\) and \(u_K\), variationally selects a unique normalizer
\(S_3\), fixes its orientation and produces the isometric sector--dodecaphase
intertwining through Reynolds averaging and
polar decomposition. The dodecaphase--Witt
step is a second incidence isometry. In the initial labeled gauge,
the two arrows do not yet realize a single representation
of \(S_3\); the joint recalibration constructed at the end of the chapter transports
the dodecad, origin, orientation, chamber and memory, and closes that
equivariance without altering the fixed-labeling no-go.
The Reynolds averages, representation decompositions and
polar factor used in this section are the standard ones
\cite{Serre1977Representations,HornJohnson2013}.

\begin{figure}[htbp]
\centering
\begin{tikzpicture}[
  font=\footnotesize,
  box/.style={draw,rounded corners=2mm,align=center,minimum width=2.35cm,minimum height=.98cm},
  arr/.style={-{Stealth[length=2.2mm]},thick},
  link/.style={<->,thick,dashed}
]
  \node[box,draw=hmtblue,fill=hmtlightblue] (seed) at (0,2.2)
    {semillas\\\(104\,976\) translated};
  \node[box,draw=hmtblue,fill=hmtlightblue] (u) at (3,2.2)
    {emissions\\\(\mathcal U_6,\ |\mathcal U_6|=468\)};
  \node[box,draw=hmtgold,fill=hmtlightgold] (ppi) at (6,2.2)
    {microfibra\\\(\Pi_\pi^{(468)},\ \operatorname{rank}=5\)};
  \node[box,draw=hmtred,fill=hmtlightred] (d) at (9,2.2)
    {positive defect\\\(D_{\rm switch}\)};
  \node[box,draw=hmtviolet,fill=hmtlightviolet] (lam) at (12,2.2)
    {capacidad residual\\\(\lambda_*=1-\delta_*\)};

  \draw[arr,hmtblue] (seed)--node[above=16pt,font=\scriptsize,fill=white,inner sep=1pt]{\(F\)} (u);
  \draw[arr,hmtgold] (u)--node[above=16pt,font=\scriptsize,fill=white,inner sep=1pt]{\(\tau_{468}\)} (ppi);
  \draw[arr,hmtred] (ppi)--node[above=16pt,font=\scriptsize,fill=white,inner sep=1pt]{\(D_{\rm switch}\)} (d);
  \draw[arr,hmtviolet] (d)--node[above=16pt,font=\scriptsize,fill=white,inner sep=1pt]{\(\tau+\mathrm{CR}\)} (lam);

  \node[box,draw=hmtgold,fill=hmtlightgold] (g) at (1.5,-.2)
    {sector oriented \(t=7\)\\\(L_{\rm or}\otimes(H_G,P_G)\)};
  \node[box,draw=hmtgreen,fill=hmtlightgreen] (v) at (5,-.2)
    {dodecaphase\\\(V_{11},P_3\)};
  \node[box,draw=hmtgreen,fill=hmtlightgreen] (e) at (8.5,-.2)
    {Witt incidence\\\(E_{30},Q_3\)};
  \node[box,draw=hmtgray,fill=hmtpaper] (t) at (12,-.2)
    {reader spectral\\\(\mathcal L_{\beta,*}\to T_{\rm band}\)};

  \draw[arr,hmtgold] (g)--node[above=16pt,font=\scriptsize,fill=white,inner sep=1pt]{\(W_{GK}\)} (v);
  \draw[arr,hmtgreen] (v)--node[above=16pt,font=\scriptsize,fill=white,inner sep=1pt]{\(U_W=I/6\)} (e);
  \draw[arr,hmtgray] (lam)--(t);
  \draw[arr,hmtred] (v) to[bend left=20]
    node[below=3pt,font=\scriptsize,fill=white,inner sep=1pt]{\((3/11)\Delta_4\)} (d);
\end{tikzpicture}

\caption{Architecture of the pre-dimensional closure. The upper row constructs
the measure and action from the emission quotient; the lower row transports
the oriented sector \(t=7\) to the dodecaphase and then the dodecaphase to
Witt. The solid arrows are proved projector isometries. The
first is fixed by the variational selection of the
\cref{thm:identificacion-isometrica-tres-proyectores}; the second is not
\(H_*\)-equivariant in the fixed gauge and becomes equivariant through the
full gauge change of the \cref{thm:recalibracion-hstar-witt}.}
\label{fig:entrelazador-accion}
\end{figure}

\section{Typing of the objects associated with \texorpdfstring{\(3/11\)}{3/11}}

Let
\[
 V_{12}=\mathbb R^{12},
 \qquad
 V_{11}=\mathbf 1^\perp
 =\left\{v\in\mathbb R^{12}:\sum_{p=1}^{12}v_p=0\right\}.
\]
The icosahedral action declared on the twelve positions decomposes as
\[
 V_{12}=V_1\oplus V_3\oplus V_{3'}\oplus V_5.
\]
Denotemos by \(P_3\) the orthogonal projector on the first summand
three-dimensional. Then
\[
 P_3^2=P_3=P_3^*,
 \qquad
 \operatorname{rank}P_3=3,
 \qquad
 \tau_{11}(P_3)
 =\frac1{11}\operatorname{Tr}_{V_{11}}P_3
 =\frac3{11}.
\]

There are three nearby objects and a fourth object of a different type:
\begin{enumerate}
  \item the point--hexad incidence transports \(P_3\) to a rank-three
  projector \(Q_3\) on a Witt module;
  \item survival in \(t=7\) distinguishes three branches among eleven and
  produces a projector \(P_G\) with normalized trace \(3/11\);
  \item an oriented intertwiner identifies
  \(P_3V_{11}\) with the oriented realization of \(P_GH_G\);
  \item a Witt star has seven even modes and four odd ones, so
  its normalized trace also equals \(3/11\), but it contains negative spectrum
  and is not a projector.
\end{enumerate}
The distinction between projector and involution will be essential to the positivity of the action law in the following section.

\section{The point--hexad incidence matrix}

Let \(\mathcal H\) be the set of the \(132\) hexads of the design
\[
 W_{12}=S(5,6,12).
\]
We define
\[
 I:\mathbb R^{12}\longrightarrow\mathbb R^{\mathcal H},
 \qquad
 (Iv)(H)=\sum_{p\in H}v_p,
\]
or, equivalently,
\[
 I_{H,p}=\mathbf 1_{\{p\in H\}}.
\]
This operator records only incidence. It contains no \(\alpha\), units,
masses, CKM, Barbero, or any physical target.

\begin{lemma}[Metric Identity of Incidence]
\label{lem:incidencia-metrica-witt}
We have
\[
 \boxed{I^*I=36I_{12}+30J_{12}.}
\]
Therefore, for every \(v\in V_{11}\),
\[
 \|Iv\|^2=36\|v\|^2.
\]
\end{lemma}

\begin{proof}
In an \(S(5,6,12)\) system, each point belongs to
\[
 r=\frac{\binom{11}{4}}{\binom54}=66
\]
hexads, and each pair of points belongs to
\[
 \lambda_2=\frac{\binom{10}{3}}{\binom43}=30
\]
hexads. The diagonal entries of \(I^*I\) have value \(66\) and the
nondiagonal entries have value \(30\). Hence
\[
 I^*I=(66-30)I_{12}+30J_{12}.
\]
Since \(J_{12}\) vanishes on \(\mathbf1^\perp\), the second
identity is obtained.
\end{proof}

To identify the spectral branch selected by \(I\), we introduce
\[
 G_{H,H'}=\binom{|H\cap H'|}{4}.
\]

\begin{lemma}[Spectral identity of incidence]
\label{lem:incidencia-espectral-witt}
With \(J_{132,12}\) denoting the rectangular all-ones matrix,
\[
 \boxed{GI=30I+15J_{132,12}.}
\]
In particular, if \(v\in V_{11}\), then
\[
 G(Iv)=30Iv.
\]
\end{lemma}

\begin{proof}
The entry \((H,p)\) of \(GI\) is
\[
 \sum_{H'\ni p}\binom{|H\cap H'|}{4}.
\]
The exhaustive census of \(W_{12}\) gives \(45\) when \(p\in H\) and \(15\)
when \(p\notin H\). Consequently, the entry equals
\(30I_{H,p}+15\). Reconstruction of the \(132\) hexads from the extended ternary
Golay code checks the \(132\cdot12\) entries as
integers.
\end{proof}

Define
\[
 E_{30}^{\rm inc}:=I(V_{11})\subseteq\ker(G-30I).
\]
Its dimension is eleven by the \cref{lem:incidencia-metrica-witt}. The calculation
modulo \(1009\) obtains
\[
 \operatorname{rank}_{\mathbb F_{1009}}(G-30I)=121.
\]
Since reduction modulo \(1009\) gives a lower bound for the rational
rank, it follows that
\(\dim_{\mathbb Q}\ker(G-30I)\le11\). The preceding inclusion provides
eleven linearly independent vectors in that kernel. Therefore,
\[
 E_{30}^{\rm inc}=\ker(G-30I).
\]
From now on we will write \(E_{30}\).

\begin{theorem}[Canonical dodecaphase entangler--Witt]
\label{thm:entrelazador-dodecafase-witt}
The operator
\[
 \boxed{
 U_W=\frac16 I\big|_{V_{11}}:V_{11}\longrightarrow E_{30}
 }
\]
is a isometry surjective. For every
\(g\in\operatorname{Aut}(W_{12})\cong M_{12}\),
\[
 U_W\rho_{12}(g)=\rho_{\mathcal H}(g)U_W.
\]
Among the equivariant point--hexad kernels, the isometry condition and
the positive sign on incidences uniquely determine \(U_W\).
\end{theorem}

\begin{proof}
The isometry follows from \(I^*I=36I\) over \(V_{11}\). Its image has
dimension eleven and is contained in \(E_{30}\), which also has dimension eleven;
it is therefore surjective. Equivariance translates
\[
 p\in H\quad\Longleftrightarrow\quad g(p)\in g(H).
\]

For uniqueness, \(M_{12}\) is transitive on the flags
\((p,H)\), with \(p\in H\), and on the antiflags. Every equivariant kernel
has two values, \(a\) on flags and \(b\) on
antiflags:
\[
 T=aI+b(J-I)=(a-b)I+bJ.
\]
The term \(J\) vanishes on \(V_{11}\), so there \(T\) is a
multiple of \(I\). The isometry forces \(|a-b|=1/6\) and the positive
normalization selects \(a-b=1/6\).
\end{proof}

\section{Positive transport of the tritic projector}

The entangler allows defining
\[
 \boxed{Q_3=U_WP_3U_W^*.}
\]
We obtain
\[
 Q_3^2=Q_3=Q_3^*,
 \qquad
 \operatorname{rank}Q_3=3,
\]
and the diagram
\[
\begin{CD}
V_{11} @>{U_W}>> E_{30}\\
@V{P_3}VV @VV{Q_3}V\\
V_{11} @>{U_W}>> E_{30}
\end{CD}
\]
commutes. In particular,
\[
 \boxed{\tau_{V_{11}}(P_3)=\tau_{E_{30}}(Q_3)=\frac3{11}.}
\]

The canonicity of \(U_W\) belongs to the labeled design and its
automorphism action. It must be distinguished from the screen \(A_5/C_5\) used to
define \(P_3\): in the constructed labeled realization, only the identity
of that concrete action of \(A_5\) preserves the fixed set of hexads. Hence
\(U_W\) canonically transports the already given \(P_3\) toward
\(Q_3\), but does not by itself convert the action \(A_5/C_5\) into a
subaction of \(M_{12}\). The bridge to the oriented sector will be constructed separately
from the marked data \(K\) and \(B_K\).

\section{Signed index of the Witt star}

Let \(B\) be a tetrad. The four hexads that contain it are written
\[
 H_i=B\sqcup P_i,\qquad |P_i|=2,
\]
where the four petals \(P_i\) partition the complement of \(B\).
The star involution \(R_B\) fixes \(B\) and exchanges the two points
of each petal. Its cycle type on twelve points is \(1^4\,2^4\). In
\(V_{11}\), and by means of \(U_W\) also in \(E_{30}\), its spectrum is
\[
 \operatorname{Spec}(R_B)=\{1^{[7]},(-1)^{[4]}\}.
\]
Therefore,
\[
 \boxed{\tau_{11}(R_B)=\frac{7-4}{11}=\frac3{11}.}
\]

\begin{theorem}[Obstruction projector--star]
\label{thm:no-go-proyector-estrella}
There is no isomorphism \(F:V_{11}\to E_{30}\) such that
\[
 FP_3=R_BF.
\]
\end{theorem}

\begin{proof}
If it existed, \(P_3\) and \(R_B\) would be similar. However,
\[
 \operatorname{Spec}(P_3)=\{1^{[3]},0^{[8]}\},
 \qquad
 \operatorname{Spec}(R_B)=\{1^{[7]},(-1)^{[4]}\}.
\]
Their minimal polynomials are \(x(x-1)\) and
\((x-1)(x+1)\), respectively. Similarity is impossible.
\end{proof}

The equality \(3/11\) thus has two readings on the same module. For
\(Q_3\) it is the mass of a positive rank-three sector. For \(R_B\) it is
a parity index. A positive action must involve \(Q_3\); the
star preserves orientation and distinguishes sheets but does not replace the
projector.

\section{homogeneous reduction \texorpdfstring{\(11\to3\to1\)}{11-3-1} in \texorpdfstring{\(t=7\)}{t=7}}

The preceding reading of \(3/11\) in \(t=8\) compared three local roots
with eleven continuations of a single root. The numerator and denominator did not
belong to the same space. The detailed table locates the correct object
in \(t=7\): there are eleven local branches; all eleven survive one horizon, three
survive two, and one survives three,
\[
 \boxed{11\longrightarrow3\longrightarrow1.}
\]

\begin{center}
\small
\begin{tabularx}{\textwidth}{cYcc}
\toprule
index&triple of words&\(h=2\)&\(h=3\)\\
\midrule
0&\texttt{200112|212221|110110}&3&0\\
1&\texttt{200121|212212|110110}&3&1\\
2&\texttt{200211|212122|110110}&3&0\\
3&\texttt{202122|210211|110110}&0&0\\
4&\texttt{202212|210121|110110}&0&0\\
5&\texttt{202221|210112|110110}&0&0\\
6&\texttt{210111|202222|110110}&0&0\\
7&\texttt{210222|202111|110110}&0&0\\
8&\texttt{212112|200221|110110}&0&0\\
9&\texttt{212121|200212|110110}&0&0\\
10&\texttt{212211|200122|110110}&0&0\\
\bottomrule
\end{tabularx}
\end{center}

Let \(\mathcal B_7\) be the set of these eleven branches and
\[
 H_G=\ell^2(\mathcal B_7).
\]
We define
\[
 P_G\delta_b=
 \begin{cases}
 \delta_b,&b\text{ admits continuation to horizon two},\\
 0,&\text{otherwise}.
 \end{cases}
\]
Then
\[
 P_G^2=P_G=P_G^*,
 \qquad
 \operatorname{rank}P_G=3,
 \qquad
 \boxed{\tau_{H_G}(P_G)=\frac3{11}.}
\]
This is the third positive realization.

\section{The residual action of \texorpdfstring{\(S_3\)}{S3} and the positive regions of \texorpdfstring{\(\pi\)}{pi}}

A permutation of the three final positions, applied simultaneously to
the first two words of each branch and leaving the third fixed, preserves
\(\mathcal B_7\). This rule defines an action of \(S_3\). Its orbits
have sizes
\[
 3+3+1+1+3=11.
\]
The support of \(P_G\), formado by the indices \(0,1,2\), is one orbit
complete, so that \(P_G\) conmuta with \(S_3\).

The character of the representation on the eleven branches, evaluated at the identity,
a transposition, and a cycle of order three, is
\[
 \chi_{\mathcal B_7}=(11,5,2)
 =5\chi_{\rm triv}+3\chi_{\rm std}.
\]
On the image of \(P_G\),
\[
 \chi_G=(3,1,0)
 =\chi_{\rm triv}+\chi_{\rm std}.
\]

The three positive regions of the microfibre of \(\pi\) have signatures
\[
 339555,\qquad393555,\qquad933555.
\]
The tails of the first word of the three surviving branches are
\[
 112,\qquad121,\qquad211.
\]
The coordinate substitution \(1\mapsto3\), \(2\mapsto9\) gives
\[
\begin{array}{c|c|c}
\text{branch}&\text{tail}&\text{positive region}\\ \hline
0&112&339555\\
1&121&393555\\
2&211&933555.
\end{array}
\]

\begin{theorem}[Horizontal--vertical bridge of the positive microfiber]
\label{thm:puente-pi-puerta-s3}
The preceding rule is an \(S_3\)-equivariant bijection between the three
positive regions of \(\pi\) and the three branches selected by \(P_G\).
The map preserves the position of the unique exceptional symbol.
\end{theorem}

\begin{proof}
Both sets are the three words obtained by placing a single exceptional
symbol in one of three positions. The substitution is applied coordinate by
coordinate and therefore commutes with every permutation of positions. The
exceptional position uniquely recovers the source element.
\end{proof}

\section{Dependence of orientation in the dodecaphase entangler--Witt}

For any normalizer \(H=N_{A_5}(C_3)\simeq S_3\) of a subgroup of
order three in \(A_5\), we define the abstract character of the quotient
\[
 \operatorname{sgn}_H:H\longrightarrow H/C_3\simeq C_2
 \longrightarrow\{+1,-1\},
 \qquad
 \operatorname{sgn}_H(h)=
 \begin{cases}
 +1,&h\in C_3,\\
 -1,&h\in H\setminus C_3.
 \end{cases}
\]
This character is not the ambient sign of a permutation of five letters:
every element of \(A_5\) has ambient sign \(+1\). It is the
intrinsic character of the reflection in the quotient \(H/C_3\). With this convention, the
character of the icosahedral sector restricted to \(H\) is
\[
 \chi_{P_3V_{11}|H}=(3,-1,0)
 =\chi_{\operatorname{sgn}_H}+\chi_{\rm std}.
\]
By contrast,
\[
 \chi_{\operatorname{im}P_G}=(3,1,0)
 =\chi_{\rm triv}+\chi_{\rm std}.
\]

\begin{corollary}[Obstruction to the unoriented braider]
\label{cor:obstruccion-entrelazador-incidencia-p3}
There is no \(S_3\)-equivariant isomorphism
\[
 \operatorname{im}P_G\longrightarrow P_3V_{11}
\]
for the non-twisted actions.
\end{corollary}

\begin{proof}
Isomorphic representations have the same character. Their values on a
transposition are \(1\) and \(-1\).
\end{proof}

Let \((e_0,e_1,e_2)\) the basis posicional of
\(\mathbb R^3_{\rm perm}\). We defines
\[
 L_{\rm or}
 =\bigwedge\nolimits^3\mathbb R^3_{\rm perm},
 \qquad
 \ell_{\rm or}=e_0\wedge e_1\wedge e_2
\]
the orientation line of the triplet, equipped with the positive anchor
\(\ell_{\rm or}\). The action of \(S_3\) on this line is precisely its
abstract sign character. Upon tensoring by it,
\[
 \chi_{L_{\rm or}\otimes\operatorname{im}P_G}
 =(3,-1,0),
\]
of mode that the character obstruction disappears. The equality of
caracteres proves the existence of one class of isomorphisms, but aun not
selects one. The selection is obtained of two witnesses already present in the
dodecaphase.

\section{Variational selection of the normalizer \texorpdfstring{\(S_3\)}{S3}}

Let
\[
 K=(234,543,140,729,659,824,621,58,914,794,146,601)
\]
the decimal chart of the seal. The common minimal witness of the two readings is
\[
 B_K=H_e\cap P_\pi
 =\{p:K_p\ge729\}
 =\{4,6,9,10\},
\]
in the mathematical numbering \(1,\ldots,12\). En the equivalent convention
with zero-based indices, the same set is \(\{3,5,8,9\}\). We define two
vectors of the sector \(P_3\):
\[
 b_K=P_3\left(\mathbf1_{B_K}-\frac{|B_K|}{12}\mathbf1\right),
 \qquad
 u_K=P_3\left(K-\overline K\,\mathbf1\right).
\]
The first satisfies exactly.
\[
 \|b_K\|^2=1.
\]

The group \(A_5\) contains ten subgroups \(C_3\). For each one, let
\[
 H(C_3)=N_{A_5}(C_3)\simeq S_3
\]
and let be the sign projector
\[
 e_{{\rm sgn},H}
 =\frac16\sum_{h\in H}\operatorname{sgn}_H(h)\rho(h).
\]
Here \(\operatorname{sgn}_H\) is the abstract character of \(H/C_3\)
defined above, not the —trivial— ambient sign of \(A_5\).
The orientation energy associated with the common witness is
\[
 \mathcal E_B(H)=\|e_{{\rm sgn},H}b_K\|^2.
\]
The exact census in \(\mathbb Q(\sqrt5)\) produces a unique maximum:
\[
 \mathcal E_B(H_*)
 =\frac23+\frac{2}{15}\sqrt5.
\]
The second level is
\[
 \frac13+\frac{2}{15}\sqrt5,
\]
so that the exact gap is
\[
 \boxed{\mathcal E_B(H_*)-\mathcal E_B(H_{\rm segundo})=\frac13.}
\]
The same \(H_*\) is the unique maximum when \(b_K\) is replaced by \(u_K\).
Therefore, the selector does not depend on a single encoding of the witness.

In this labeling,
\[
 H_*=N_{A_5}\bigl(\langle c_*\rangle\bigr),
\qquad
 c_*=(2\,3\,4).
\]
The rule is natural under simultaneous relabeling by \(A_5\): when
\(B_K\), \(K\), and \(P_3\) are transported, the maximum is conjugated by the same
element.

\section{Selection of orientation and presentation of \texorpdfstring{\(S_3\)}{S3}}

Within of \(H_*\), consider
\[
 s(g)=\langle b_K,\rho(g)u_K\rangle.
\]
The maximum among the two order-three elements uniquely selects
\[
 c_*=(0,1,3,4,2),
\]
in image notation of the permutation of five points, and the maximum among
the three involutions selects
\[
 r_*=(1,0,4,3,2).
\]
It is verified exactly.
\[
 r_*c_*r_*=c_*^{-1}.
\]
The presentation of the oriented sector by the cycle of positions
\(c=(0\,1\,2)\) and the reflection \(r=(0\,2)\), which fixes the central position,
therefore determines a unique isomorphism
\[
 \theta:S_3^{\rm or}\xrightarrow{\sim}H_*,
 \qquad
 \theta(c)=c_*,
 \qquad
 \theta(r)=r_*.
\]
Here \(S_3^{\rm or}\) denotes the group of relabelings of the oriented sector
\(t=7\), generated by \(c\) and \(r\); it introduces no additional operatorial
class.
The only branch that reaches horizon three is precisely the central branch
\(1\); this datum fixes which of the three positions remains fixed under
reflection.

\section{Reynolds average and polar factor}

Let
\[
 W_{\rm or}=L_{\rm or}\otimes\mathbb R^3_{\rm perm}
\]
the oriented sector triple. With the anchor already fixed, we take
\[
 \boxed{w=\ell_{\rm or}\otimes e_1,}
\]
where \(e_1\) represents the central position, the only one reaching horizon
three. Define
\[
 A
 =\sum_{g\in S_3}
 \rho_{W_{\rm or}}(g)\,
 |w\rangle\langle b_K|\,
 \rho_{V_{11}}(\theta(g))^{-1}.
\]
By construction,
\[
 A\rho_{V_{11}}(\theta(g))
 =\rho_{W_{\rm or}}(g)A,
\qquad
 A=AP_3.
\]
The exact computation in \(\mathbb Q(\sqrt5)\) obtains
\[
 \operatorname{rank}A=3.
\]
Moreover, \(G_A=AA^*\) has all its diagonal entries equal to
\[
 3+\frac25\sqrt5
\]
and all its off-diagonal entries equal to
\[
 \frac52+\frac35\sqrt5.
\]
Its two isotype eigenvalues are
\[
 \lambda_{\rm sgn}=8+\frac85\sqrt5,
 \qquad
 \lambda_{\rm std}=\frac12-\frac15\sqrt5.
\]
Both are strictly positive. In particular, \(G_A^{-1/2}\) exists, and
the polar decomposition
\[
 \boxed{U_{GK}=G_A^{-1/2}A}
\]
is one equivariant isometry of \(P_3V_{11}\) on \(W_{\rm or}\):
\[
 U_{GK}^*U_{GK}=P_3,
 \qquad
 U_{GK}U_{GK}^*=I_{W_{\rm or}}.
\]

Let
\[
 J_G:W_{\rm or}\longrightarrow L_{\rm or}\otimes H_G
\]
the isometric inclusion in the three selected branches. Then
\[
 J_GJ_G^*=I_{L_{\rm or}}\otimes P_G.
\]
The sector--dodecaphase partial isometry is denoted
\[
 \boxed{
 W_{GK}=U_{GK}^*J_G^*
 :L_{\rm or}\otimes H_G\longrightarrow V_{11}
 }
\]
satisfies
\[
 W_{GK}^*W_{GK}
 =I_{L_{\rm or}}\otimes P_G,
\qquad
 W_{GK}W_{GK}^*=P_3,
\]
and, therefore,
\[
 P_3W_{GK}
 =W_{GK}
 =W_{GK}(I_{L_{\rm or}}\otimes P_G).
\]

\begin{theorem}[Relative Isometric Identification of Three Projectors]
\label{thm:identificacion-isometrica-tres-proyectores}
Relative to the action \(A_5/C_5\), the projector \(P_3\), the seal \(K\), the doubly constructed witness \(B_K\), the oriented finite sector \(t=7\), the abstract character \(\operatorname{sgn}_{H_*}\), the anchor \(\ell_{\rm or}\), and the variational rule \(\mathcal E_B\), there exists a unique normalized partial isometry \(W_{GK}\) that identifies \(P_G\) with \(P_3\). Independently, the incidence isometry \(U_W\) identifies \(P_3\) with \(Q_3\). Consequently, the following chain of isometric identifications of projectors is obtained:
\[
\begin{CD}
L_{\rm or}\otimes H_G
 @>{W_{GK}}>>
 V_{11}
 @>{U_W}>>
 E_{30}\\
@V{I_{L_{\rm or}}\otimes P_G}VV
 @VV{P_3}V
 @VV{Q_3}V\\
L_{\rm or}\otimes H_G
 @>{W_{GK}}>>
 V_{11}
 @>{U_W}>>
 E_{30}.
\end{CD}
\]
Typed traces satisfy
\[
 \tau_{L_{\rm or}\otimes H_G}(I_{L_{\rm or}}\otimes P_G)
 =\tau_{V_{11}}(P_3)
 =\tau_{E_{30}}(Q_3)
 =\frac3{11}.
\]
The statement, in the fixed gauge, isometrically identifies three
projectors; it does not yet assert that their three spaces are the same
representation of a common group.
\end{theorem}

\begin{proof}
The energy \(\mathcal E_B\) selects a unique normalizer \(H_*\);
the scoring criterion \(s(g)\) selects a unique pair \((c_*,r_*)\), and with it a
unique isomorphism \(\theta\). The Reynolds average is then
equivariant for the actions of \(S_3^{\rm or}\) and \(H_*\) related
by \(\theta\), and has rank three. The positive part of its polar
decomposition fixes the normalization without introducing a basis. The preceding
partial identities prove that \(W_{GK}\) identifies
\(P_G\) and \(P_3\). On the other hand, the
\cref{thm:entrelazador-dodecafase-witt} identifies \(P_3\) and \(Q_3\)
as projectors through \(U_W\). These two operator identities make
that fixed-gauge diagram commute, without extending the equivariance of the
first arrow to the second there.
\end{proof}

\begin{proposition}[Obstruction to triple equivariance in the fixed gauge]
\label{prop:no-go-equivarianza-hstar-witt}
The preceding construction does not prove that \(U_W\) is
\(H_*\)-equivariant. Under the declared labeling, the fixed action of
\(A_5/C_5\) on the twelve points preserves the set of hexads of
\(W_{12}\) only for the identity element. Therefore, the
\(H_*\)-equivariant intertwining of \(W_{GK}\) does not extend, in that
gauge, to a common representation on \(E_{30}\).
\end{proposition}

\begin{proof}
The first assertion separates equivariance domains:
\(W_{GK}\) uses \(H_*\subset A_5\), whereas the incidence theorem
proves equivariance of \(U_W\) for
\(\operatorname{Aut}(W_{12})\). The exact reproducible census applies the sixty
elements of action \(A_5/C_5\) to the labeled hexad system; only the
identity preserves that system. Since \(H_*\) contains six elements, its
action cannot enter by restriction of that fixed action in
\(\operatorname{Aut}(W_{12})\).
\end{proof}

The preceding proposition isolates exactly the fixed gauge; it is not an
abstract obstruction to transporting the calibration jointly. The
governing owner of \(S_3\)--Witt recalibration constructs the equivariant monomorphism and
second arrow as follows.

\begin{theorem}[Joint \texorpdfstring{\(H_*\)}{H*}--Witt recalibration]
\label{thm:recalibracion-hstar-witt}
Let
\(\bar\rho_{\rm src}:H_*\to S_{12}\) be the point action on the twelve cosets
determined by \(c_*\), \(r_*\) and the isomorphism \(\theta\) of the
oriented sector, and let \(\rho_{\rm src}\) be its linearization on \(V_{11}\).
In the one-based image convention, the permutation
\[
 \boxed{\sigma=(1,5,2,6,7,4,3,10,11,9,8,12)}
\]
defines the monomorphism
\[
 \boxed{
 \iota_\sigma:H_*\hookrightarrow\operatorname{Aut}(W_{12})\cong M_{12},
 \qquad
 \iota_\sigma(h)=
 \sigma\,\bar\rho_{\rm src}(h)\,\sigma^{-1}.
 }
\]
If the calibration is transported in full through
\[
 W_{12}^{\sigma}=\{\sigma^{-1}B:B\in W_{12}\},
 \qquad
 \widetilde U_W:=U_WP_\sigma,
\]
where \(P_\sigma\) is the permutation unitary, then, under the
corresponding identification of the transported incidence space with
\(E_{30}\), for every \(h\in H_*\),
\[
 \boxed{
 \widetilde U_W\,\rho_{\rm src}(h)
 =\rho_{E_{30}}\!\bigl(\iota_\sigma(h)\bigr)\,\widetilde U_W.
 }
\]
In particular,
\[
 \widetilde Q_3=\widetilde U_WP_3\widetilde U_W^*
\]
is the Witt realization of the same projector and the chain
\[
 I_{L_{\rm or}}\otimes P_G
 \xleftrightarrow{\ W_{GK}\ }
 P_3
 \xleftrightarrow{\ \widetilde U_W\ }
 \widetilde Q_3
\]
carries, after identifying the source action through \(\theta\), a
common representation of \(H_*\), with the target action given by
\(\iota_\sigma\). Recalibration jointly transports
\(C_W,W_{12},I\), the dodecad, origin, orientation, chamber and
memory; it does not consist in replacing \(U_W\) in isolation.
\end{theorem}

\begin{proof}
The exact certificate reconstructs the ternary Golay code, the \(132\)
hexads, the group \(M_{12}\) of order \(95040\), the source \(S_3\) and a
semiregular target \(S_3\), both of order six. It checks for each
\(h\in H_*\) that
\(
\sigma\bar\rho_{\rm src}(h)\sigma^{-1}\in M_{12}
\), that \(W_{12}^{\sigma}\) is \(H_*\)-invariant and that its point--hexad
incidence is equivariant. Injectivity of \(\iota_\sigma\) follows
because it is conjugation of the faithful source action. Moreover,
\[
 P_\sigma\rho_{\rm src}(h)
 =\rho_{12}\!\bigl(\iota_\sigma(h)\bigr)P_\sigma.
\]
Composing this identity with the \(M_{12}\) equivariance of \(U_W\)
yields the boxed identity. Since \(P_3\) commutes with \(H_*\), its
transport \(\widetilde Q_3\) commutes with \(\iota_\sigma(H_*)\); the
already proved equivariance of \(W_{GK}\) completes the composition. The same
certificate verifies that the original design is not invariant under the source
action, so the fixed-gauge no-go is preserved without contradiction.
\end{proof}

\textsc{Material owner.} This result is governed by theorem
\textsc{HMT--RDF--EXC--S3}, sections 3.1--3.3 of the governing Witt--Moonshine exceptional
propagation delta, and its reproducible certificate
\texttt{verificar\_recalibracion\_s3\_witt.py}, whose output is
\texttt{PASS\_RECALIBRACION\_S3\_WITT}. The status is
\texttt{FORMALIZACION\_NUEVA/CERTIFICADO\_NUEVO}; uniqueness of the
gauge change is not postulated.

Before using \(B_K\), \(u_K\), orientation and the polar factor, there were
\(10\cdot6\cdot4=240\) labeled maps, or \(40\) after identifying
internal relabelings of the sector. The preceding rules reduce that
set to one. This is the verifiable content of the word
\emph{canonical} for \(W_{GK}\), within the declared rules. The common
action on the Witt module comes from the explicit change \(\sigma\); it is not
asserted that this gauge change is unique or canonical.

The variational rule \(\mathcal E_B\) uses only prior HMT data
and fixes the class in which uniqueness is obtained. It does not identify the star
involution with \(P_3\), since their spectra remain distinct.

\section{Comparison of the readings of \texorpdfstring{\(t=7\) and \(t=8\)}{t=7 and t=8}}

In \(t=8\) there are three local roots and eleven extensions to horizon nine of
the unique root that continues. The sizes of the three fibers are
\[
 (11,0,0).
\]
The natural branch--chain matrix has one row of eleven ones and two
zero rows, so its rank is one. In \(t=9\), moreover, the visible quotient
changes from \(3/11\) to \(3/6\). That fraction was not the trace of a
stable projector. The correction does not eliminate the sector: it shifts the reading
to the homogeneous object \(P_G\) of \(t=7\), where rank and dimension
belong to the same space.

\section{Consequence for a positive action}

The law of action may u now a unique class isometric of projectors
positive in tris spaces of support:
\[
 P_G
 \xleftrightarrow{\ W_{GK}\ }
 P_3
 \xleftrightarrow{\ U_W\ }
 Q_3.
\]
The star \(R_{B_K}\) remains reserved for parity and orientation. The
projector \(P_G\) records survival; \(P_3\) records dodecaphase
polarization; \(Q_3\) records its incidence realization in Witt. This
separation of functions is the mathematical premise of the measure and of the positive
law in the next section. In the fixed gauge the chain expresses
isometric equivalence of projectors without a common representation. The
\cref{thm:recalibracion-hstar-witt} transports the complete calibration and
turns the chain with \(\widetilde U_W\) and \(\widetilde Q_3\) into a
common representation of \(H_*\), without identifying the signed star with the
positive projector.
